% sections/5_conclusions.tex
\section{Conclusiones}

El presente informe ha demostrado con éxito la resolución numérica de la ecuación de Poisson empleando diversos enfoques computacionales acelerados.

En primer lugar, el análisis del sistema cilíndrico ha validado el algoritmo de Gauss-Seidel en coordenadas cilíndricas ($r, z$), reproduciendo el decaimiento logarítmico del potencial radial predicho teóricamente para la sección media longitudinal. Pequeñas desviaciones analíticas son atribuibles a efectos de borde en la geometría finita.

En segundo lugar, se ha presentado un avanzado sistema de trabajo automatizado que procesa imágenes geométricas arbitrarias mediante visión artificial con OpenCV (\texttt{cv2}). Se ha comprobado que la detección de bordes y la identificación de contornos permiten mapear eficazmente geometrías no euclídeas a condiciones de contorno físicas válidas para la ecuación de Poisson, ya sean líneas de potencial o densidades de carga. El estudio de la tolerancia iterativa ha revelado la importancia de alcanzar grados de convergencia profundos (tolerancia $\leq 10^{-3}$) para obtener isolíneas de potencial suaves y estables.

Finalmente, el contraste con el ejercicio académico 9.1 ha corroborado que el método basado en visión artificial puede reproducir de forma trivial y exacta los resultados obtenidos mediante el enfoque estándar cuando se interpreta un dibujo conceptual sombreado como densidades de carga proporcionales. Este hecho ratifica la utilidad y flexibilidad del sistema de visión artificial desarrollado para resolver problemas de electrostática en presencia de distribuciones de carga complejas definidas geométricamente.

